\section{Infrastructure and Data-Cleaning Details}
\label{app:infra-cleaning}

This appendix documents the engineering substrate used to produce the reported benchmark results. The goal is not to prescribe one implementation stack, but to make clear which parts of \arena{} are part of the benchmark contract and which parts are deployment choices. We describe the implementation paths that matter for reproducibility: how a run is launched, how wrappers are constrained, how intermediate results are recovered, and how the dataset pool is cleaned before benchmark distillation.

\subsection{Layered execution architecture}

\arena{} separates algorithm integration from benchmark execution through four layers.

\begin{center}
\small
\begin{tabular}{p{0.20\linewidth}p{0.31\linewidth}p{0.39\linewidth}}
\toprule
\textbf{Layer} & \textbf{Main artifact} & \textbf{Responsibility} \\
\midrule
Configuration & \texttt{toolbox\_config.json}, \texttt{envs\_config.json} & Map each tool to a wrapper class and environment \\
Front-end API & \texttt{SymbolicRegressor} & Expose \texttt{fit}, \texttt{predict}, equation extraction, and artifact export \\
Isolation runner & \texttt{subprocess\_runner.py} & Execute wrapper commands in the configured environment \\
Benchmark runner & \texttt{benchmarks/runner.py} & Load data, inject contracts, evaluate metrics, and write artifacts \\
\bottomrule
\end{tabular}
\end{center}

The benchmark never calls an upstream algorithm repository directly. Instead, a run is issued as a JSON command containing the action, data array, seed, budget, and framework parameters. The subprocess runner imports the appropriate wrapper module, executes the action, serializes the fitted wrapper state, and returns a JSON response. This design lets the benchmark keep a stable interface even when the underlying algorithms expose incompatible APIs.

\subsection{Wrapper contract}

Each algorithm wrapper implements a minimal estimator interface:
\begin{verbatim}
fit(X, y)
predict(X)
get_optimal_equation()
get_total_equations()
export_canonical_symbolic_program()
\end{verbatim}
The first four methods are the execution interface. The last method is the symbolic-artifact interface: it exports a \texttt{CanonicalSymbolicProgram} containing the raw equation, normalized expression, variables, fitted constants when available, operator set, AST size, tree depth, and validation notes. When a wrapper cannot directly provide a complete canonical artifact, the result layer applies tool-specific normalizers to translate outputs such as PySR expressions, gplearn prefix trees, DSO/uDSR expressions, PyOperon strings, TPSR/E2ESR candidates, iMCTS outputs, and LLM-generated Python functions into the shared representation.

The benchmark runner also injects an explicit data contract into every wrapper:
\begin{verbatim}
n_features
feature_names
target_name
seed
timeout_in_seconds
exp_path
exp_name
\end{verbatim}
Wrappers validate that the declared feature count and names match the actual training matrix. This avoids a common SR-integration failure mode: a method silently using the wrong number of input variables or emitting an out-of-range variable token. For LLM-based methods, \arena{} separates the real data contract from the prompt contract. The runner can expose canonical prompt variables while preserving original feature names, target names, metadata-derived background, and optional anonymization switches.

\subsection{Environment isolation and why Docker is not the primary boundary}

\arena{} uses per-tool conda environments plus subprocess isolation as the canonical evaluation boundary. The current 12-method panel requires incompatible Python and package stacks: for example, the base GP/PySR/PyOperon environment uses Python 3.10, DSO/uDSR and E2ESR use Python 3.7, TPSR uses Python 3.9, and LLM-backed methods use a separate LLM environment. Some methods also require external assets such as pretrained weights, Julia/PySR caches, GPU libraries, or API-backed LLM configuration.

Docker is useful for lightweight demonstration and single-environment deployment, but it is not used as the primary benchmark boundary for three practical reasons. First, a single container that includes all 12 algorithms would be large, slow to rebuild, and brittle because it would have to mix mutually incompatible Python and native dependencies. Second, the reported experiments run on heterogeneous CPU/GPU cluster nodes where user-level conda environments, mounted datasets, local model weights, Julia caches, and private LLM configuration files are easier to manage than privileged container builds. Third, the benchmark requires fine-grained process supervision: per-run timeout recovery, stdout/stderr forwarding, result-file recovery, and minute-level progress extraction are naturally handled by the subprocess runner. The artifact therefore records environment specifications and wrapper versions, while Docker remains an optional packaging layer rather than the source of evaluation semantics.

\subsection{Result, timeout, and progress semantics}

Every run writes a canonical \texttt{result.json} containing the tool name, dataset identity, seed, sanitized parameters, status, termination reason, final expression, canonical symbolic artifact, train/validation/ID/OOD metrics, runtime, and experiment directory. Sensitive keys such as API keys or tokens are removed from saved parameter dictionaries.

\arena{} distinguishes engineering failure from algorithmic failure. A run can be:
\begin{itemize}[leftmargin=*]
    \item \texttt{ok}: completed or budget-exhausted with a recovered, evaluable best-so-far expression;
    \item \texttt{timed\_out}: budget exhausted without a recoverable evaluable expression;
    \item \texttt{no\_valid\_output}: the method ran but produced no expression that can be evaluated on the required splits;
    \item \texttt{error}: wrapper, environment, data-loading, or execution error.
\end{itemize}
For snapshot-capable tools, the runner polls the algorithm's current-best state every 60 seconds and writes a minute-level JSON file under \texttt{progress/}. If the final expression is numerically unstable or a timeout occurs, the runner can recover the latest evaluable snapshot and mark the run as budget-exhausted with output instead of discarding the run.

\subsection{Dataset sources and acquisition}

The dataset repository is maintained separately from the paper source and the algorithm code. In the released artifact, datasets are accessed through stable manifest entries rather than copied into the experiment directory. Each manifest row stores the dataset identity, source family, subgroup, relative data path, split files, metadata file, and ground-truth formula file when available. This keeps the benchmark reproducible while avoiding multiple drifting copies of the same raw data.

The canonical data root is organized as
\begin{verbatim}
sim-datasets-data/<family>/<subgroup>/<dataset>/
\end{verbatim}
or, for compact classical families,
\begin{verbatim}
sim-datasets-data/<family>/<dataset>/
\end{verbatim}
The same layout is used locally, on the compute cluster, and in the public dataset package. The dataset package can be obtained either as a Git-LFS backed repository from Hugging Face or ModelScope, or through the lightweight \texttt{sim-datasets} Python package. During cluster execution, run manifests resolve paths such as \texttt{sim-datasets-data/...} against the configured dataset root, not against a partial checkout inside the code repository.

\begin{center}
\small
\begin{tabular}{p{0.22\linewidth}p{0.28\linewidth}p{0.40\linewidth}}
\toprule
\textbf{Source family} & \textbf{Role in the reservoir} & \textbf{Normalization notes} \\
\midrule
SRSD & Formula-backed scientific discovery tasks with Feynman-style equations and dummy-variable variants & Preserve SRSD subgroup labels, dummy-variable status, source citations, and \texttt{formula.py}; normalize Feynman identifiers for duplicate detection \\
LLM-SRBench & Physics, chemistry, materials, biological growth, and transformed equation-discovery tasks & Preserve semantic feature/target descriptions for prompt-controlled LLM evaluation while also supporting anonymized-variable runs \\
SRBench 1.0 & Feynman and Strogatz-style benchmark tasks used in prior SR evaluation & Convert source naming into the shared family/subgroup/basename scheme and audit cross-source Feynman duplicates \\
SRBench2025 first-principles & Compact first-principles tasks with small sample counts and explicit scientific formulas & Preserve original descriptions, units when available, and source-derived or table-derived OOD splits \\
Nguyen, Keijzer, Korns, Vladislavleva & Classical synthetic SR families used as controlled baselines & Generate deterministic train/validation/ID/OOD splits from formula specifications and record the generation rule in metadata \\
\bottomrule
\end{tabular}
\end{center}

The resulting \reservoir{} is therefore not a single upstream benchmark. It is a normalized union of formula-backed tasks from SRSD, LLM-SRBench, SRBench 1.0, SRBench2025 first-principles tasks, and classical synthetic families. Source family and subgroup are retained as first-class metadata fields so later selection stages can control coverage rather than treating the reservoir as an anonymous pool.

\subsection{Dataset format and validation}

All datasets are normalized into one executable directory format:
\begin{verbatim}
train.csv
valid.csv
id_test.csv
ood_test.csv        # optional but required for OOD analysis
metadata.yaml
formula.py          # required for ground-truth reservoir tasks
\end{verbatim}
The CSV files contain only numeric columns. The target column is not inferred from position during evaluation; it is read from \texttt{metadata.yaml}. The ordered feature list in metadata is the authoritative input signature and must match the non-target CSV columns. This matters for SR because variable order is part of the equation semantics: a valid expression over the wrong feature ordering is treated as an integration error rather than an algorithmic result.

The metadata records dataset name, source family, subgroup, split files and sample counts, ordered feature names, target name, feature and target descriptions when available, train/OOD ranges, license/source information, upstream resources, and the ground-truth formula file. When a source benchmark provides physical or semantic feature descriptions, \arena{} preserves them in metadata. LLM-based methods can then be evaluated either with semantic background enabled or with anonymized variables, while non-LLM methods receive the same numeric arrays and data contract.

Split construction follows a deterministic policy. If the upstream source already provides train, validation, ID-test, or OOD-test splits, those splits are preserved after header and type normalization. If a source provides only a single formula-backed table, \arena{} uses a fixed random seed to produce train/validation/ID splits and, when a reliable formula and feature domain are available, synthesizes an OOD split by sampling outside the training support and replaying the ground-truth formula. If no reliable OOD definition is available, the pipeline does not fabricate an OOD split; the gap is recorded in metadata and the dataset is excluded from OOD-dependent analyses.

For ground-truth tasks, \texttt{formula.py} is part of the benchmark contract. It exposes a deterministic callable evaluated using the metadata feature order. The validation script checks that required splits exist, files are not Git-LFS pointer stubs, CSV headers agree across splits, target columns exist, declared sample counts match row counts, all feature and target columns are numeric, \texttt{metadata.yaml} agrees with the CSV schema, and \texttt{formula.py} reproduces split targets within numeric tolerance. It also verifies that formula replay produces finite values on the declared OOD split; otherwise the dataset is marked for repair or manual review before it can enter reservoir selection.

\subsection{Cleaning and audit pipeline}

The data-cleaning pipeline is deliberately split into mechanical repairs and human-review gaps. Mechanically derivable fields are filled only when they can be computed from local files: missing \texttt{ground\_truth\_formula.file} is filled when \texttt{formula.py} exists, split ratios are computed from sample counts, formula NMSE is computed by replaying \texttt{formula.py} on split CSVs, and target OOD ranges are computed from \texttt{ood\_test.csv}. Fields requiring semantic judgment, such as citation text or feature descriptions for black-box variables, are not hallucinated; they are emitted into a manual-gap report.

\arena{} also audits common dataset-packaging failures before selection. Git-LFS pointer files are detected before interpreting CSV headers. All selected manifests are resolved against the real dataset root rather than a partial repository mirror. Dataset identity fields are normalized so that local paths, remote paths, and \texttt{sim-datasets-data/...} relative paths map to the same identity. This prevents false missing-file reports and avoids accidentally evaluating a result on the wrong dataset.

\subsection{Semantic duplicate removal and final materialization}

Before constructing the final benchmark, \arena{} performs semantic duplicate detection on the internal 100-task clean pool used during development. For Feynman-style tasks, names are normalized into a canonical Feynman identifier. For general formula-backed tasks, \texttt{formula.py} is parsed, function arguments are normalized to \texttt{x0}, \texttt{x1}, \ldots, and a symbolic expression hash is computed from the simplified expression tree. If multiple tasks share a semantic key, the pipeline retains the highest-priority representative and fills the removed slot from \textsc{Candidate-200} using nearest-neighbor replacement under family, subgroup, selection-mode, advantage-side, feature-count, sample-count, and formula-complexity similarity.

The final materialization step intentionally excludes result-derived columns such as priority score, probe gap, quality score, stability score, selected advantage side, and selection mode. It uses only static attributes: family, subgroup, dummy-variable flag, feature count, sample scale, and formula complexity. A randomized search and audit report family/subgroup balance and static numeric summaries. Finally, selected datasets are materialized through manifest files and stable links, so the benchmark is auditable without duplicating raw data directories.
